\section{Derivatives of Regular Expressions with Positive Lookahead}
\label{sec:derivatives-positive-lookahead}
% ============================================================

This section extends the Brzozowski derivative construction recalled in Section~\ref{sec:derivatives-of-regular-expressions} to $\mathrm{REwPLA}$. We first define the left quotient of string pairs and of languages with string pair introduced in Section~\ref{sec:languages-with-prefix-contexts}. We then define derivatives for $\mathrm{REwPLA}$ and prove that, for every expression in $\mathrm{REwPLA}$, the language denoted by its derivative coincides with the left quotient of its language.
% Since positive lookahead is
% interpreted using languages of string pairs rather than ordinary
% languages, the residual objects of the derivative operation are elements
% of $\mathcal C_\Sigma=\mathcal P(\Sigma^*\times\Sigma^*)$.

% A string pair $(u,v)$ records the two components introduced in
% Section~\ref{sec:languages-with-prefix-contexts}: the first component
% $u$ is the remaining main string, while the second component $v$ is the
% remaining right-context string. The derivative operation consumes the
% first component first. Only after the main component has been exhausted
% does it consume the second component.


% Their semantic
% counterparts on string pairs and on languages of string pairs are denoted
% by $\partial_a^{\mathsf P}$ and $\partial_w^{\mathsf P}$.

% ------------------------------------------------------------
\subsection{Left Quotient of String Pairs and Languages of String Pairs}
\label{sec:stringpair-derivatives}
% ------------------------------------------------------------

For a symbol $a\in\Sigma$, we define the left quotient of a string pair by
% $\partial_a^{\mathsf P}:\Sigma^*\times\Sigma^*\to\mathcal C_\Sigma$ by
\begin{equation}
a^{-1}(u,v)
=
\begin{cases}
\{(a^{-1}u,v)\},
& \quad\text{if }a\preceq u,\\
\{(\varepsilon,a^{-1}v)\},
& \quad \text{if }u=\varepsilon
  \text{ and }a\preceq v,\\
\emptyset,
& \quad \text{otherwise}.
\end{cases}
\label{eq:stringpair-symbol-derivative}
\end{equation}
Here, $a^{-1}u$ and $a^{-1}v$ denote the left quotients of strings by symbol $a$, as defined in
Section~\ref{sec:derivatives-of-regular-expressions}. The left quotient of a string pair by a string is defined inductively by
$
\varepsilon^{-1}(u,v)=\{(u,v)\},
\:\:\:
(aw)^{-1}(u,v)
=
\bigcup_{\xi\in a^{-1}(u,v)}
w^{-1}(\xi).
\label{eq:stringpair-string-derivative}
$
In particular, for identity $\mathbf{1}$, 
$
a^{-1}(\varepsilon,\varepsilon) = \emptyset .
\label{eq:stringpair-fully-discharged}
$ For example,
\begin{equation}
\begin{aligned}
a^{-1}(ab,cd)
&=\{(b,cd)\}, \quad
&
ab^{-1}(ab,cd)
&=\{(\varepsilon,cd)\},
\\
abc^{-1}(ab,cd)
&=\{(\varepsilon,d)\}, \quad
&
abcd^{-1}(ab,cd)
&=\{(\varepsilon,\varepsilon)\}.
\end{aligned}
\label{eq:stringpair-derivative-example}
\end{equation}
% Thus, reading $ab$ exhausts the main component, while reading $cd$
% subsequently discharges the right-context component. The derivative
% operation does not project a string pair to an ordinary string.

% ------------------------------------------------------------
% \subsection{Derivatives of Languages of String Pairs}
% \label{sec:stringpair-language-derivatives}
% ------------------------------------------------------------

The left quotient is lifted pointwise from string pairs to
languages of string pairs. For $R\in\mathcal C_\Sigma$ and
$a\in\Sigma$, define
$
a^{-1}R
=
\bigcup_{\xi\in R}
a^{-1}(\xi).$
\label{eq:stringpair-language-symbol-derivative}
The left quotient of a language of string pair by a string is defined inductively by $
\varepsilon^{-1}R=R,
\:
(aw)^{-1}R
=
w^{-1}
\bigl(a^{-1}R\bigr).
\label{eq:stringpair-language-word-derivative}
$
Equivalently,
$
w^{-1}R
=
\bigcup_{\xi\in R}
w^{-1}(\xi).
\label{eq:stringpair-language-word-derivative-pointwise}
$ Clearly, $w^{-1}R\in\mathcal{C}_{\Sigma}$.

\subsection{Derivatives of \textbf{$\mathrm{REwPLA}$}}
We define a derivative operation on $\mathrm{REwPLA}$ that maps an expression to another expression in $\mathrm{REwPLA}$. Following the Brzozowski notation introduced in Section~\ref{sec:derivatives-of-regular-expressions}, we write $D_a(r)$ and $D_w(r)$ for the derivatives of an expression $r\in\mathrm{REwPLA}$ with respect to a symbol $a\in\Sigma$ and a string $w\in\Sigma^*$, respectively. 

To facilitate the definition of derivatives for $\mathrm{REwPLA}$, we introduce two Boolean functions. $r \in \mathrm{REwPLA},$ function $\lambda_1$ indicates whether $\mathcal{M}(r)$ contains a pair whose first component is empty string (i.e., $\varepsilon$), while $\lambda$ indicates whether the identity pair belongs to $\mathcal{M}(r)$. Formally,
$
\lambda_1(r)=\mathsf{T}
\:\Longleftrightarrow\:
\exists \xi\in\mathcal{M}(r),\ \pi_1(\xi)=\varepsilon,
$
and
$
\lambda(r)=\mathsf{T}
\:\Longleftrightarrow\:
(\varepsilon,\varepsilon)\in\mathcal{M}(r),
$
with both functions returning $\mathsf{F}$ otherwise.

We introduce an auxiliary derivative $D'_a(r)$, called the \emph{virtual derivative}, which captures derivations that consume the symbol $a$ through a string constraint associated with an empty main-string component. Thus, $D'_a(r)$ is relevant only when $\lambda_1(r)=\mathsf{T}$. For $a,b\in\Sigma$, the virtual derivative is defined inductively as follows:
\begingroup
\footnotesize
\begin{equation}
\begin{gathered}
D'_a(\emptyset)=\emptyset,
\qquad
D'_a(\varepsilon)=\emptyset,
\qquad
D'_a(b)=\emptyset,
\quad
D'_a(r^*)=
\begin{cases}
D'_a(r)r^*,
& \lambda_1(r)=\mathsf{T},\\
\emptyset,
& \text{otherwise},
\end{cases}
\quad
D'_a(\mathsf{LA}(r))
=
\mathsf{LA}\bigl(D_a(r)\bigr),
\\
D'_a(r+s)=
\begin{cases}
D'_a(r)+D'_a(s),
& \lambda_1(r)\land\lambda_1(s),\\
D'_a(r),
& \lambda_1(r)=\mathsf{T}\land\lambda_1(s)=\mathsf{F},\\
D'_a(s),
& \lambda_1(r)=\mathsf{F}\land\lambda_1(s)=\mathsf{T},\\
\emptyset,
& \text{otherwise},
\end{cases}
\quad
D'_a(rs)=
\begin{cases}
\emptyset,
& \lambda_1(r)=\mathsf{F}\lor\lambda_1(s)=\mathsf{F},\\
D'_a(r)D'_a(s),
& \lambda(r)=\mathsf{F}\land\lambda(s)=\mathsf{F},\\
D'_a(r)D'_a(s)+D'_a(s),
& \lambda(r)=\mathsf{T}\land\lambda(s)=\mathsf{F},\\
D'_a(r)D'_a(s)+D'_a(r),
& \lambda(r)=\mathsf{F}\land\lambda(s)=\mathsf{T},\\
D'_a(r)D'_a(s)+D'_a(r)+D'_a(s),
& \lambda(r)=\mathsf{T}\land\lambda(s)=\mathsf{T}.
\end{cases}
\end{gathered}
\label{eq:rewpla-virtual-derivative}
\end{equation}
\endgroup
The derivative $D_a(r)$ of an expression $r\in\mathrm{REwPLA}$ with respect to a symbol $a\in\Sigma$ captures both real and virtual derivations while preserving the associated string constraints. It is defined inductively as follows:
\begingroup
\footnotesize
\begin{equation}
\begin{gathered}
D_a(\emptyset)=\emptyset,
\qquad
D_a(\varepsilon)=\emptyset,
\qquad
D_a(b)=
\begin{cases}
\varepsilon, & \text{if } a=b,\\
\emptyset, & \text{otherwise},
\end{cases}
\quad
D_a(r+s)=D_a(r)+D_a(s),
\qquad
D_a(r^*)=D_a(r)r^*,
\\[1mm]
D_a(\mathsf{LA}(r))
=
\mathsf{LA}\bigl(D_a(r)\bigr),
\quad
D_a(rs)=
\begin{cases}
D_a(r)s+D'_a(r)D_a(s),
& \lambda_1(r)=\mathsf{T}\land\lambda(r)=\mathsf{F},
\\[1mm]
D_a(r)s+D'_a(r)D_a(s)+D_a(s),
& \lambda(r)=\mathsf{T},
\\[1mm]
D_a(r)s,
& \text{otherwise}.
\end{cases}
\end{gathered}
\label{eq:rewpla-symbol-derivative}
\end{equation}
\endgroup
Here, $D'_a(r)$ denotes the virtual component of the derivative, whereas $D_a(r)$ represents the complete derivative. In particular, $D'_a(r)$ accounts for the case in which the main-string component is empty and the symbol $a$ is consumed through the associated string constraint.

\begin{example}
Consider the expression $e=(\mathsf{LA}(ab)+a)(a+b)^*,$
and compute its derivative with respect to $a$.
\begin{align*}
D_a(e)
&=D_a\bigl(\mathsf{LA}(ab)+a\bigr)(a+b)^*
  +D'_a\bigl(\mathsf{LA}(ab)+a\bigr)
   D_a\bigl((a+b)^*\bigr)
\\
&=\bigl(D_a(\mathsf{LA}(ab))+D_a(a)\bigr)(a+b)^*
  +D'_a(\mathsf{LA}(ab))
   D_a(a+b)(a+b)^*
\\
&=\bigl(\mathsf{LA}(D_a(ab))+\varepsilon\bigr)(a+b)^*
  +\mathsf{LA}(D_a(ab))
   \bigl(D_a(a)+D_a(b)\bigr)(a+b)^*
\\
&=\bigl(\mathsf{LA}(b)+\varepsilon\bigr)(a+b)^*
  +\mathsf{LA}(b)
   (\varepsilon+\emptyset)(a+b)^*
\\
&=\bigl(\mathsf{LA}(b)+\varepsilon\bigr)(a+b)^*
  +\mathsf{LA}(b)(a+b)^*
\\
&=(a+b)^*+\mathsf{LA}(b)(a+b)^* .
\end{align*}

\end{example}

We define the nullability function
\[
\lambda:\mathrm{REwPLA}\longrightarrow\{\emptyset,\varepsilon\}
\]
by
\[
\lambda(r)=
\begin{cases}
\varepsilon,
& \text{if }(\varepsilon,\varepsilon)\in\mathcal{M}(r),\\
\emptyset,
& \text{otherwise}.
\end{cases}
\]
Thus, $\lambda(r)=\varepsilon$ exactly when $r$ admits the identity
pair, i.e., when both the main-string component and the string-constraint
component are empty.

\begin{theorem}[Closure under derivatives]
For every expression $r\in\mathrm{REwPLA}$, every symbol
$a\in\Sigma$, and every word $w\in\Sigma^*$,
\[
D_a(r)\in\mathrm{REwPLA}
\qquad\text{and}\qquad
D_w(r)\in\mathrm{REwPLA}.
\]
Moreover, the auxiliary virtual derivative $D'_a(r)$ occurring in the
definition of $D_a(r)$ is also an expression in $\mathrm{REwPLA}$.
\end{theorem}

\begin{proof}
We first prove simultaneously, by structural induction on $r$, that
$D_a(r)$ and $D'_a(r)$ belong to $\mathrm{REwPLA}$.

For the base expressions $\emptyset$, $\varepsilon$, and $b\in\Sigma$,
the claim follows immediately from the definitions, since their
derivatives are either $\emptyset$ or $\varepsilon$.

For $r+s$, the derivatives are constructed from the derivatives of
$r$ and $s$ using only union. For $rs$, every branch in the definitions
of $D_a(rs)$ and $D'_a(rs)$ is constructed from derivatives of $r$ and
$s$ using union and concatenation. Hence the claim follows from the
induction hypothesis.

For $r^*$, the derivatives are constructed from derivatives of $r$
using concatenation and Kleene star. For $\mathsf{LA}(r)$, both
$D_a(\mathsf{LA}(r))$ and $D'_a(\mathsf{LA}(r))$ are constructed by
applying $\mathsf{LA}$ to $D_a(r)$. Therefore, all resulting expressions
belong to $\mathrm{REwPLA}$.

Thus $D_a(r)\in\mathrm{REwPLA}$ for every symbol $a$. Since derivatives
with respect to words are defined by
\[
D_\varepsilon(r)=r,
\qquad
D_{aw}(r)=D_w(D_a(r)),
\]
an induction on the length of $w$ yields
$D_w(r)\in\mathrm{REwPLA}$.
\end{proof}

\begin{theorem}[Derivative equation for ordinary regular expressions]
For every ordinary regular expression $r\in\mathrm{RE}$,
\[
r
\equiv
\lambda(r)
+
\sum_{a\in\Sigma} aD_a(r).
\]
Equivalently,
\[
\mathcal{L}(r)
=
\mathcal{L}(\lambda(r))
\cup
\bigcup_{a\in\Sigma}
a\,\mathcal{L}(D_a(r)).
\]
\end{theorem}

\begin{proof}
Let $w\in\mathcal{L}(r)$. If $w=\varepsilon$, then
$\lambda(r)=\varepsilon$, and hence
$w\in\mathcal{L}(\lambda(r))$.

Otherwise, $w=ax$ for some $a\in\Sigma$ and $x\in\Sigma^*$. By the
correctness of derivatives,
\[
x\in\mathcal{L}(D_a(r))
\quad\Longleftrightarrow\quad
ax\in\mathcal{L}(r).
\]
Hence $w=ax\in a\mathcal{L}(D_a(r))$.

Conversely, $\varepsilon\in\mathcal{L}(\lambda(r))$ only when
$\varepsilon\in\mathcal{L}(r)$. Moreover, if
$x\in\mathcal{L}(D_a(r))$, derivative correctness gives
$ax\in\mathcal{L}(r)$. Therefore both sides denote the same language.
\end{proof}